\documentclass[12pt]{article}
\usepackage{amsmath, amsthm, amssymb}
\usepackage[margin=1in]{geometry}
 
\newtheorem*{problem}{Problem}
\newtheorem*{theorem}{Theorem}
 
\theoremstyle{remark}
\newtheorem*{remark}{Remark}
 
\title{Main Idea: Composite Map Range Bounded by Minimum of Individual Ranges}
\author{Liam Abrams}
\date{April 25, 2026}
 
\begin{document}
\maketitle
 
%------------------------------------------------------------------
\section*{Chapter 3 -- Problem 23 of LADR4E by Axler}
%------------------------------------------------------------------
 
\begin{problem}
Suppose $U$ and $V$ are finite-dimensional vector spaces and $S \in \mathcal{L}(V, W)$
and $T \in \mathcal{L}(U, V)$. Prove that
\[
    \dim \operatorname{range} ST \leq
    \min\{\dim \operatorname{range} S,\; \dim \operatorname{range} T\}.
\]
\end{problem}
 
\begin{proof}
We establish each inequality separately.
 
\medskip
\noindent\textbf{Claim 1: $\dim \operatorname{range} ST \leq \dim \operatorname{range} S$.}
 
\medskip
For any $u \in U$, we have $ST(u) = S(Tu)$, which is an element of
$\operatorname{range} S$ by definition. Therefore every vector in
$\operatorname{range} ST$ belongs to $\operatorname{range} S$, i.e.,
\[
    \operatorname{range} ST \subseteq \operatorname{range} S.
\]
Since a subspace cannot have larger dimension than any subspace that contains it:
\[
    \dim \operatorname{range} ST \leq \dim \operatorname{range} S.
\]
 
\medskip
\noindent\textbf{Claim 2: $\dim \operatorname{range} ST \leq \dim \operatorname{range} T$.}
 
\medskip
The composition $ST$ first applies $T: U \to V$, then applies $S$. Every output
of $ST$ is of the form $S(v)$ for some $v \in \operatorname{range} T$, so
\[
    \operatorname{range} ST = \{S(v) : v \in \operatorname{range} T\}.
\]
In other words, $\operatorname{range} ST$ is the image of $\operatorname{range} T$
under $S$. Viewing the restriction $S|_{\operatorname{range} T}$ as a linear map
from $\operatorname{range} T$ to $W$ and applying the Fundamental Theorem of
Linear Maps:
\[
    \dim(\operatorname{range} T)
    = \dim\bigl(\operatorname{null} S|_{\operatorname{range} T}\bigr)
      + \dim\bigl(\operatorname{range} S|_{\operatorname{range} T}\bigr).
\]
Since $\operatorname{range} ST = \operatorname{range} S|_{\operatorname{range} T}$
and dimensions are non-negative:
\[
    \dim \operatorname{range} ST
    = \dim(\operatorname{range} T)
      - \dim\bigl(\operatorname{null} S|_{\operatorname{range} T}\bigr)
    \leq \dim \operatorname{range} T.
\]
 
\medskip
\noindent\textbf{Conclusion.}

\medskip
Since $\dim \operatorname{range} ST \leq \dim \operatorname{range} S$ and
$\dim \operatorname{range} ST \leq \dim \operatorname{range} T$, the quantity
$\dim \operatorname{range} ST$ is a lower bound for both $\dim \operatorname{range} S$
and $\dim \operatorname{range} T$, and therefore cannot exceed their minimum:
\[
    \dim \operatorname{range} ST
    \leq \min\{\dim \operatorname{range} S,\; \dim \operatorname{range} T\}. \qedhere
\]

\end{document}